\section{Related Work}

\textbf{Physarum shortest path.} The slime mould resolves a maze into the
shortest connecting path~\cite{nakagaki2000maze}; Tero et al.\ formalised
this as adapting conductivities driven by a Poiseuille-flow pressure
solve~\cite{tero2007mathematical}, later used to reproduce Tokyo rail-like
networks~\cite{tero2010rules}. Bonifaci et al.\ proved convergence to the
shortest path for \emph{undirected} graphs~\cite{bonifaci2012physarum}. The
directedness restriction is explicit in the literature: Zhang et al.\ note
the model is designed only for undirected networks and propose a directed
formulation~\cite{zhang2014directed}, which we do not implement. Our results
therefore characterise the canonical model on a directed network, an acknowledged
limitation, rather than reporting a defect.

\textbf{The prior three-way comparison.} Combining Physarum, ACO and
Dijkstra is not new. Zhang et al.\ benchmark exactly these
three~\cite{zhang2014improved}, reporting that their Physarum variant and
Dijkstra both attain optimal paths ``with one hundred percent'' while ACO
degrades with network size, on Erd\H{o}s--R\'enyi graphs of 15--2{,}000
nodes with static weights, no congestion and no disruption. Two points
matter. Their method needs \emph{no} extraction machinery, since the path is read
straight off the flux field, which separates cleanly on such graphs, so we
cannot claim that Physarum evaluations generally conflate provenance;
the honest contrast is that extraction is a formality there and decisive
here. And their ACO reaches zero optimal-path accuracy beyond 1{,}000 nodes,
whereas ours attains 43--44\% exact match on a larger and harder network,
which establishes our ACO as a strong baseline by the standard of the prior
art.

\textbf{Extraction and termination.} Wang et al.\ define a dominant-path
procedure procedurally identical to our maximum-flux
walk~\cite{wang2021transition}, and state that at convergence it yields
exactly the shortest path. On our network it yields a directed-legal path on
43.0\% of solves and the optimum on 23.3\%. The results are compatible: their
examples are complete graphs and small-worlds of mean degree 6--600, assumed
undirected, where the walk \emph{cannot} dead-end. Our network has mean
out-degree 2.46 with pervasive one-ways.

\textbf{ACO on traffic networks.} Zhang et al.\ propose orientation
guidance weighting edges by angular alignment with the
destination~\cite{zhang2018novel}, evaluated on one 528-node network with
static travel times against conventional ACO only; our goal-directed
ablation is related but distinct and we extend rather than contradict it.
Tuan et al.\ report an improved ACO beating conventional ACO, Dijkstra and a
genetic algorithm~\cite{tuan2026improved}, but evaluate on METR-LA under
fluctuating traffic~\cite{li2018diffusion} without measuring recovery after a
closure.

\textbf{Position.} A systematic pass over the open citation databases,
including complete forward-citation sweeps of
both~\cite{zhang2014improved} and~\cite{watanabe2011traffic}, returns no
study combining the three families on a real road network under
simulator-generated congestion with closure recovery as a measured outcome;
no work pairing Physarum with SUMO was found at all. The contribution
claimed here is the substrate and the measurement, not the three-algorithm
combination, which~\cite{zhang2014improved} already reported.
